\documentclass[12pt]{article}
\usepackage[utf8]{inputenc}
\usepackage[T1]{fontenc}
\usepackage[margin=1in]{geometry}
\usepackage{setspace}
\usepackage{xcolor}
\usepackage[colorlinks=true, urlcolor=blue!60!black, linkcolor=blue!60!black]{hyperref}
\setlength{\parindent}{2em}
\setlength{\parskip}{1.25em}
\renewcommand{\baselinestretch}{1.0}
\begin{document}
\begin{center}
{\LARGE \textit{Florida Prepaid Postsecondary Education Expense Board v. College Savings Bank}, 527 U.S. 627 (1999)\par}
\vspace{0.8em}
{\large Case Brief\par}
\end{center}
\vspace{1em}

\section{Memory Jogger}

Applying \textit{City of Boerne}'s ``congruence and proportionality'' test, the Court struck down the Patent Remedy Act's abrogation of state sovereign immunity under \S~5 of the Fourteenth Amendment because Congress compiled virtually no legislative record of widespread, unremedied, nonnegligent state patent infringements violating the Due Process Clause (\href{https://www.courtlistener.com/opinion/118316/florida-prepaid-postsecondary-education-expense-board-v-college-savings/}{CourtListener}; \href{https://www.law.cornell.edu/supct/html/98-531.ZO.html}{LII Majority}).

\section{Detailed Case Facts}

Beginning in 1987, College Savings Bank, a New Jersey-chartered savings bank located in Princeton, marketed and sold certificates of deposit called ``CollegeSure CDs,'' which functioned as annuity contracts designed to guarantee investors sufficient funds to cover future college tuition costs (\href{https://www.law.cornell.edu/supct/html/98-531.ZO.html}{LII Majority}). College Savings obtained U.S. Patent No. 4,722,055 (the ``'055 patent'') covering its financing methodology and data-processing apparatus (\href{https://law.justia.com/cases/federal/district-courts/FSupp/948/400/2098999/}{D.N.J. Opinion}; \href{https://law.justia.com/cases/federal/district-courts/FSupp/919/756/1581067/}{Justia}).

The Florida Prepaid Postsecondary Education Expense Board (``Florida Prepaid''), an arm of the State of Florida created under Fla. Stat. \S~240.551(1), administered competing tuition prepayment contracts for Florida residents and their children (\href{https://www.law.cornell.edu/supct/html/98-531.ZO.html}{LII Majority}). College Savings alleged that Florida Prepaid directly, contributorily, and willfully infringed the '055 patent in administering its prepayment program (\href{https://www.law.cornell.edu/supct/html/98-531.ZO.html}{LII Majority}; \href{https://law.justia.com/cases/federal/district-courts/FSupp/948/400/2098999/}{D.N.J. Opinion}).

Prior to 1992, the Patent Act provided only that ``whoever'' infringed a patent was liable, 35 U.S.C. \S~271(a), which the Federal Circuit held in \textit{Chew v. California}, 893 F.2d 331 (Fed. Cir. 1990), lacked the ``unmistakably clear'' statement required by \textit{Atascadero State Hospital v. Scanlon}, 473 U.S. 234 (1985), to abrogate state sovereign immunity (\href{https://www.law.cornell.edu/supct/html/98-531.ZO.html}{LII Majority}). In response, Congress enacted the Patent and Plant Variety Protection Remedy Clarification Act of 1992 (``Patent Remedy Act''), Pub. L. No. 102-560, 106 Stat. 4230, codified at 35 U.S.C. \S\S~271(h) and 296(a), explicitly defining ``whoever'' to include any State or state instrumentality and declaring that States ``shall not be immune, under the eleventh amendment of the Constitution of the United States or under any other doctrine of sovereign immunity, from suit in Federal court by any person'' for patent infringement (\href{https://www.law.cornell.edu/supct/html/98-531.ZO.html}{LII Majority}). Congress invoked three constitutional bases: the Patent Clause (Art.~I, \S~8, cl.~8), the Interstate Commerce Clause (Art.~I, \S~8, cl.~3), and \S~5 of the Fourteenth Amendment (\href{https://www.law.cornell.edu/supct/html/98-531.ZO.html}{LII Majority}).

\section{Procedural History}

On November 7, 1994, College Savings Bank filed a patent infringement action against Florida Prepaid under 35 U.S.C. \S~271 in the U.S. District Court for the District of New Jersey, seeking declaratory and injunctive relief, damages, attorney's fees, and costs (and separately filed a Lanham Act false-advertising suit resolved in the companion case, \textit{College Savings Bank v. Florida Prepaid Postsecondary Education Expense Board}, 527 U.S. 666 (1999)) (\href{https://www.law.cornell.edu/supct/html/98-531.ZO.html}{LII Majority}; \href{https://law.justia.com/cases/federal/district-courts/FSupp/948/400/2098999/}{D.N.J. Opinion}).

After the Supreme Court held in \textit{Seminole Tribe of Florida v. Florida}, 517 U.S. 44 (1996), that Congress cannot abrogate state sovereign immunity using its Article~I powers, Florida Prepaid moved to dismiss on Eleventh Amendment grounds (\href{https://www.law.cornell.edu/supct/html/98-531.ZO.html}{LII Majority}). College Savings argued that Congress validly abrogated immunity under \S~5 of the Fourteenth Amendment to enforce the Due Process Clause, and the United States intervened to defend the statute (\href{https://www.law.cornell.edu/supct/html/98-531.ZO.html}{LII Majority}). On December 13, 1996, the District Court denied Florida Prepaid's motion to dismiss the patent claim, 948 F. Supp. 400 (D.N.J. 1996) (\href{https://law.justia.com/cases/federal/district-courts/FSupp/948/400/2098999/}{D.N.J. Opinion}). On June 30, 1998, the U.S. Court of Appeals for the Federal Circuit affirmed, 148 F.3d 1343 (Fed. Cir. 1998), holding that patents are Fourteenth Amendment property and that the Patent Remedy Act was a proportionate enforcement measure under \textit{City of Boerne v. Flores}, 521 U.S. 507 (1997) (\href{https://law.justia.com/cases/federal/fed-cir/1998/148-f3d-1343.html}{Fed. Cir. Opinion}). The Supreme Court granted certiorari, heard oral argument on April 20, 1999, and reversed and remanded on June 23, 1999 (\href{https://www.oyez.org/cases/1998/98-531}{Oyez}; \href{https://www.law.cornell.edu/supct/html/98-531.ZS.html}{LII Syllabus}).

\section{Judicial Votes}

The Supreme Court decided the case by a \textbf{5--4} vote, with no concurring opinions (\href{https://www.law.cornell.edu/supct/html/98-531.ZS.html}{LII Syllabus}; \href{https://www.oyez.org/cases/1998/98-531}{Oyez}):

\begin{itemize}
\item \textbf{Majority (5):} Chief Justice William H. Rehnquist (author of the opinion of the Court), joined by Associate Justices Sandra Day O'Connor, Antonin Scalia, Anthony M. Kennedy, and Clarence Thomas (\href{https://www.law.cornell.edu/supct/html/98-531.ZS.html}{LII Syllabus}).
\item \textbf{Dissenting (4):} Associate Justice John Paul Stevens (author of the dissenting opinion), joined by Associate Justices David H. Souter, Ruth Bader Ginsburg, and Stephen G. Breyer (\href{https://www.law.cornell.edu/supct/html/98-531.ZS.html}{LII Syllabus}; \href{https://www.law.cornell.edu/supct/html/98-531.ZD.html}{LII Dissent}).
\end{itemize}

\section{Holding}

\paragraph{Judgment.}
Judgment of the U.S. Court of Appeals for the Federal Circuit \textbf{reversed and remanded} (\href{https://www.courtlistener.com/opinion/118316/florida-prepaid-postsecondary-education-expense-board-v-college-savings/}{CourtListener}; \href{https://www.law.cornell.edu/supct/html/98-531.ZS.html}{LII Syllabus}).

\paragraph{Rule Announced.}
Although Congress unequivocally expressed its intent to abrogate state sovereign immunity in the Patent Remedy Act and patents are ``property'' protected by the Fourteenth Amendment's Due Process Clause, Congress may not abrogate state sovereign immunity under Article~I (\textit{Seminole Tribe}), and the Patent Remedy Act cannot be sustained as ``appropriate'' enforcement legislation under \S~5 of the Fourteenth Amendment (\href{https://www.law.cornell.edu/supct/html/98-531.ZO.html}{LII Majority}). To qualify as valid remedial or prophylactic \S~5 legislation under \textit{City of Boerne v. Flores}, Congress must identify a pattern of conduct transgressing the Fourteenth Amendment's substantive provisions---here, intentional or reckless state patent infringement coupled with absent or constitutionally inadequate state remedies---and tailor its legislative scheme so that there is ``congruence and proportionality'' between the constitutional injury and the means adopted (\href{https://www.law.cornell.edu/supct/html/98-531.ZO.html}{LII Majority}).

\paragraph{Holding vs. Dicta.}
The \textbf{holding} is that 35 U.S.C. \S\S~271(h) and 296(a) fail \textit{City of Boerne}'s congruence-and-proportionality test and therefore do not validly abrogate States' Eleventh Amendment immunity from federal patent-infringement suits (\href{https://www.law.cornell.edu/supct/html/98-531.ZS.html}{LII Syllabus}). Key subsidiary rulings integral to the holding include: (1) patents are ``property'' under the Due Process Clause; (2) state patent infringement violates due process only if it is nonnegligent and the State provides no adequate remedy; and (3) the Patent Remedy Act's indiscriminate, nationwide abrogation lacked legislative-record support and tailoring (\href{https://www.law.cornell.edu/supct/html/98-531.ZO.html}{LII Majority}). By contrast, Chief Justice Rehnquist's discussion of hypothetical narrowing features that Congress \textit{could} have adopted---such as limiting abrogation to willful infringement, infringement authorized by state policy, or States lacking adequate state-court remedies---is instructive \textbf{dicta}, as is the footnote observing that Florida provided state-law remedies via a legislative claims bill or state takings/conversion suit (\href{https://www.law.cornell.edu/supct/html/98-531.ZO.html}{LII Majority}). Similarly, the Court expressly declined to reach whether the Fifth Amendment's Just Compensation Clause (incorporated through the Fourteenth Amendment) could support abrogation because Congress never invoked it (\href{https://www.law.cornell.edu/supct/html/98-531.ZO.html}{LII Majority}).

\section{Analysis of Opinions}

\subsection{Majority Opinion (Chief Justice William H. Rehnquist)}

Chief Justice Rehnquist began with the two-part abrogation inquiry from \textit{Seminole Tribe}: (1) whether Congress unequivocally expressed its intent to abrogate state sovereign immunity, and (2) whether Congress acted pursuant to a valid grant of constitutional power (\href{https://www.law.cornell.edu/supct/html/98-531.ZO.html}{LII Majority}). Having noted that the companion case foreclosed implied constructive waiver under \textit{Parden}, Rehnquist readily found Congress's abrogation intent ``unmistakably clear'' in 35 U.S.C. \S~296(a) (\href{https://www.law.cornell.edu/supct/html/98-531.ZO.html}{LII Majority}). Under \textit{Seminole Tribe}, however, neither the Patent Clause nor the Commerce Clause in Article~I permits abrogation, leaving \S~5 of the Fourteenth Amendment (per \textit{Fitzpatrick v. Bitzer}, 427 U.S. 445 (1976)) as the sole viable constitutional basis (\href{https://www.law.cornell.edu/supct/html/98-531.ZO.html}{LII Majority}).

Applying \textit{City of Boerne v. Flores}, 521 U.S. 507 (1997), Rehnquist emphasized that \S~5 grants Congress a ``remedial'' power to enforce---not substantively redefine---the Fourteenth Amendment, requiring ``congruence and proportionality between the injury to be prevented or remedied and the means adopted to that end'' (\href{https://www.law.cornell.edu/supct/html/98-531.ZO.html}{LII Majority}). First, Rehnquist examined the legislative record and found virtually no pattern of state patent infringement: the House Report cited only two examples, the Federal Circuit found only eight suits against States in 110 years (1880--1990), and witnesses and the bill's sponsor (Rep. Kastenmeier) conceded that States generally complied with patent law (\href{https://www.law.cornell.edu/supct/html/98-531.ZO.html}{LII Majority}).

Second, although Rehnquist agreed that patents are ``property'' under the Due Process Clause, he stressed that state patent infringement does not \textit{by itself} violate the Fourteenth Amendment. Under \textit{Zinermon v. Burch}, 494 U.S. 113 (1990), \textit{Parratt v. Taylor}, 451 U.S. 527 (1981), and \textit{Hudson v. Palmer}, 468 U.S. 517 (1984), a procedural due process violation occurs only when the State provides \textit{no} remedy or constitutionally \textit{inadequate} remedies (\href{https://www.law.cornell.edu/supct/html/98-531.ZO.html}{LII Majority}). Yet Congress barely examined state remedies; witnesses complained mostly about inconvenience and loss of national uniformity---concerns belonging to Article~I's Patent Clause, not Fourteenth Amendment due process (\href{https://www.law.cornell.edu/supct/html/98-531.ZO.html}{LII Majority}). Moreover, under \textit{Daniels v. Williams}, 474 U.S. 327 (1986), mere negligence does not ``deprive'' a person of property under the Due Process Clause, yet patent infringement under \S~271(a) is a strict-liability tort and the legislative record indicated that most state infringement was innocent or negligent (\href{https://www.law.cornell.edu/supct/html/98-531.ZO.html}{LII Majority}).

Finally, Rehnquist held the Patent Remedy Act grossly disproportionate: Congress subjected all 50 States to immediate federal-court liability for every form of direct, induced, or contributory infringement, for an indefinite duration, without limiting coverage to willful infringement, state-policy infringement, or States lacking adequate remedies (\href{https://www.law.cornell.edu/supct/html/98-531.ZO.html}{LII Majority}). Because the Act's true goals---national uniformity and placing States on equal footing with private parties---were Article~I objectives rather than Fourteenth Amendment enforcement, the abrogation was invalid (\href{https://www.law.cornell.edu/supct/html/98-531.ZO.html}{LII Majority}).

\subsection{Dissenting Opinion (Justice John Paul Stevens)}

Justice Stevens, joined by Justices Souter, Ginsburg, and Breyer, argued that the Patent Remedy Act was ``a paradigm of an appropriate exercise of Congress' \S~5 power'' (\href{https://www.law.cornell.edu/supct/html/98-531.ZD.html}{LII Dissent}). First, Stevens stressed the constitutional and statutory imperative of national uniformity in patent law: since 1800, Congress has vested \textit{exclusive} jurisdiction over patent infringement in federal courts (28 U.S.C. \S~1338(a)) and consolidated patent appeals in the Federal Circuit (\href{https://www.law.cornell.edu/supct/html/98-531.ZD.html}{LII Dissent}). Because federal law preempted state patent jurisdiction, Congress reasonably assumed state remedies were nonexistent or inadequate, and forcing patentees into 50 disparate state courts---which lack patent expertise, cannot be reviewed by the Federal Circuit, and may not even waive state immunity---undermines both uniformity and impartial due process (\href{https://www.law.cornell.edu/supct/html/98-531.ZD.html}{LII Dissent}).

Second, Stevens faulted the majority for applying \textit{Daniels v. Williams} to dismiss the statute on its face even though College Savings explicitly alleged \textit{willful} infringement (\href{https://www.law.cornell.edu/supct/html/98-531.ZD.html}{LII Dissent}). He also called it ``quite unfair'' to fault the 1992 Congress for not compiling a 50-state evidentiary record to satisfy a \textit{City of Boerne} standard the Court did not announce until five years later, when Congress in 1992 was simply meeting \textit{Atascadero}'s clear-statement rule while \textit{Pennsylvania v. Union Gas Co.}, 491 U.S. 1 (1989), remained good law (\href{https://www.law.cornell.edu/supct/html/98-531.ZD.html}{LII Dissent}).

Third, Stevens distinguished \textit{City of Boerne}. Unlike RFRA---which attempted to overturn the Court's substantive First Amendment doctrine and displaced countless state laws---the Patent Remedy Act did not alter any substantive rule of state law; it merely removed a sovereign-immunity defense in a narrow, exclusively federal category of cases (\href{https://www.law.cornell.edu/supct/html/98-531.ZD.html}{LII Dissent}). Its congruence was self-scaling: if state infringement is rare, the Act rarely operates; if state commercial patent activity grows (noting that all 50 States own patents and public universities collect hundreds of millions in royalties), the Act expands in lockstep (\href{https://www.law.cornell.edu/supct/html/98-531.ZD.html}{LII Dissent}). Stevens closed by reiterating his fundamental dissent from \textit{Seminole Tribe}'s ``judge-made doctrine of sovereign immunity'' (\href{https://www.law.cornell.edu/supct/html/98-531.ZD.html}{LII Dissent}).

\section{Examples: Future Applications}

\begin{enumerate}
\item \textbf{Same-Side (Invalid Copyright Abrogation):} Congress enacts the Copyright Remedy Clarification Act (CRCA) to abrogate state sovereign immunity for all federal copyright infringement claims, relying on a legislative record showing only a dozen or so anecdotal state infringements and failing to restrict liability to intentional infringements or States lacking adequate state remedies. Under \textit{Florida Prepaid}, the abrogation is invalid under \S~5 because it fails \textit{City of Boerne}'s congruence-and-proportionality test (directly mirroring \textit{Allen v. Cooper}, 589 U.S. 248 (2020)) (\href{https://www.oyez.org/cases/2019/18-877}{Oyez}).

\item \textbf{Same-Side (Invalid Age-Discrimination Abrogation):} Congress amends the Age Discrimination in Employment Act (ADEA) to allow state employees to sue States for money damages in federal court, without documenting a widespread pattern of unconstitutional, irrational age discrimination by state employers. Because age classifications receive only rational-basis review under the Equal Protection Clause and Congress prohibited far more conduct than the Fourteenth Amendment forbids without a supporting legislative record, the abrogation is invalid under \textit{Florida Prepaid} and \textit{City of Boerne} (mirroring \textit{Kimel v. Florida Board of Regents}, 528 U.S. 62 (2000)) (\href{https://supreme.justia.com/cases/federal/us/528/62/}{Justia}).

\item \textbf{Opposite-Side (Valid Abrogation for Heightened-Scrutiny Gender Discrimination):} Congress enacts the Family and Medical Leave Act's family-care leave provision, expressly abrogating state sovereign immunity after compiling an extensive legislative record showing pervasive, stereotype-driven gender discrimination by state employers in granting family leave, and carefully tailors the remedy (unpaid leave only, minimum tenure requirements, exclusion of high-ranking state officials). Unlike \textit{Florida Prepaid}, this abrogation is valid under \S~5 because gender discrimination triggers heightened scrutiny and Congress's narrowly tailored remedy is congruent and proportional to a documented constitutional evil (mirroring \textit{Nevada Department of Human Resources v. Hibbs}, 538 U.S. 721 (2003)) (\href{https://www.law.cornell.edu/supct/html/01-1368.ZS.html}{LII}).

\item \textbf{Opposite-Side (Targeted Statutory Abrogation for Willful Patent Takings Without State Remedies):} Suppose Congress reenacts a narrower patent-abrogation statute that applies \textit{only} when (a) a State commits \textit{willful} or \textit{intentional} patent infringement pursuant to official state policy, and (b) the patentee demonstrates that the State provides no adequate state-court judicial remedy for compensation. Under \textit{Florida Prepaid}'s own roadmap, this tightly tailored statute would likely be upheld under \S~5 because it targets actual Fourteenth Amendment due process violations rather than innocent infringement or States with adequate remedies (\href{https://www.law.cornell.edu/supct/html/98-531.ZO.html}{LII Majority}).

\item \textbf{Fence-Sitter (Conditioning Federal Research Grants on Sovereign-Immunity Waiver):} Rather than relying on \S~5 abrogation, Congress invokes its Article~I Spending Clause power to condition a state university's receipt of federal biomedical research grants (or its right to obtain federal patents from the USPTO) on the State's express waiver of Eleventh Amendment immunity in federal patent-infringement suits. This is a classic fence-sitter: while \textit{Florida Prepaid} and \textit{College Savings Bank} foreclosed unilateral abrogation and ``constructive'' waiver, \textit{South Dakota v. Dole} permits Congress to induce voluntary waivers via Spending Clause conditions---so long as the condition is unambiguous, germane to the federal funds, and not unconstitutionally coercive (a boundary heavily contested in post-\textit{Florida Prepaid} scholarship) (\href{https://scholarship.law.duke.edu/cgi/viewcontent.cgi?article=1142&context=dlj}{Duke Law Journal}).
\end{enumerate}

\section{Critique}

\paragraph{Originalist and Textualist Perspectives.}
Defenders aligned with structural federalism praise \textit{Florida Prepaid} for enforcing enumerated-powers limits and preventing Congress from using \S~5 as an end-run around \textit{Seminole Tribe} to federalize ordinary Article~I regulatory disputes (\href{https://www.repository.law.indiana.edu/cgi/viewcontent.cgi?article=1843&context=ilj}{Indiana Law Journal}; \href{https://digitalcommons.law.villanova.edu/cgi/viewcontent.cgi?article=3073&context=vlr}{Villanova Law Review}). Yet prominent textualist and originalist scholars---including John F. Manning and others examining the Eleventh Amendment's precise text---have sharply criticized the foundation on which \textit{Florida Prepaid} rests: the Eleventh Amendment's text bars only diversity suits against a State by citizens of another State or a foreign state, not federal-question suits, meaning \textit{Hans v. Louisiana} and \textit{Seminole Tribe} rely on nontextual ``presuppositions'' rather than constitutional text (\href{https://scholarship.law.uc.edu/cgi/viewcontent.cgi?article=1475&context=uclr}{University of Cincinnati Law Review}; \href{https://www.uclalawreview.org/wp-content/uploads/2019/09/47_50UCLALRev10312002-2003.pdf}{UCLA Law Review}). Other critics note that \textit{Florida Prepaid} inverted traditional facial-challenge rules (invalidating the Act even as applied to College Savings's allegation of \textit{willful} infringement) and imposed an ahistorical, quasi-administrative ``legislative record'' requirement on Congress (\href{https://repository.law.umich.edu/cgi/viewcontent.cgi?article=1794&context=mlr}{Michigan Law Review}).

\paragraph{Progressive and Doctrinal Perspectives.}
Progressive constitutional scholars---including Erwin Chemerinsky, Robert Post, and Reva Siegel---condemn \textit{Florida Prepaid} as a centerpiece of the Rehnquist Court's ``juricentric Constitution,'' which treats the Supreme Court as the sole expositor of constitutional meaning and strips Congress of the broad, \textit{McCulloch}-style discretion the Reconstruction Framers intended in \S~5 (\href{https://law.yale.edu/sites/default/files/documents/pdf/Faculty/Siegel_ProtectingTheConstitutionFromThePeople.pdf}{Yale Law School / Indiana Law Journal}; \href{http://www.stanfordlawreview.org/wp-content/uploads/sites/3/2010/04/chemerinsky.pdf}{Stanford Law Review}). Doctrinal commentators also emphasize a glaring practical asymmetry: States actively obtain and enforce hundreds of federal patents against private citizens in federal court, yet \textit{Florida Prepaid} shields States from reciprocal damages liability when they infringe private patents, undermining national patent uniformity (\href{https://verdict.justia.com/2020/04/10/how-allen-v-cooper-breaks-important-new-if-dubious-ground-on-stare-decisis}{Justia Verdict}; \href{https://repository.uclawsf.edu/cgi/viewcontent.cgi?article=3736&context=hastings_law_journal}{Hastings Law Journal}).

\paragraph{Competing Underlying Values.}
\textit{Florida Prepaid} reflects a deep collision among three constitutional values: \textbf{state sovereign dignity and dual federalism} (shielding state treasuries from private federal-court damages suits), \textbf{separation of powers and judicial supremacy} (policing the line between judicial interpretation of \S~1 and congressional enforcement under \S~5), and \textbf{remedial accountability and national uniformity} (ensuring that federally created property rights carry effective, uniform remedies even when the infringer is a state government).

\section{Key Quotations}

\begin{quote}
``We thus held that for Congress to invoke \S~5, it must identify conduct transgressing the Fourteenth Amendment's substantive provisions, and must tailor its legislative scheme to remedying or preventing such conduct.''
\end{quote}
\noindent --- Chief Justice William H. Rehnquist, Opinion of the Court, 527 U.S. at 639 (\href{https://www.courtlistener.com/opinion/118316/florida-prepaid-postsecondary-education-expense-board-v-college-savings/}{CourtListener}; \href{https://www.law.cornell.edu/supct/html/98-531.ZO.html}{LII Majority}).

\begin{quote}
``Thus, under the plain terms of the Clause and the clear import of our precedent, a State's infringement of a patent, though interfering with a patent owner's right to exclude others, does not by itself violate the Constitution. Instead, only where the State provides no remedy, or only inadequate remedies, to injured patent owners for its infringement of their patent could a deprivation of property without due process result.''
\end{quote}
\noindent --- Chief Justice William H. Rehnquist, Opinion of the Court, 527 U.S. at 643 (\href{https://www.courtlistener.com/opinion/118316/florida-prepaid-postsecondary-education-expense-board-v-college-savings/}{CourtListener}; \href{https://www.law.cornell.edu/supct/html/98-531.ZO.html}{LII Majority}).

\begin{quote}
``The need for uniformity in the construction of patent law is undoubtedly important, but that is a factor which belongs to the Article I patent-power calculus, rather than to any determination of whether a state plea of sovereign immunity deprives a patentee of property without due process of law.''
\end{quote}
\noindent --- Chief Justice William H. Rehnquist, Opinion of the Court, 527 U.S. at 645 (\href{https://www.courtlistener.com/opinion/118316/florida-prepaid-postsecondary-education-expense-board-v-college-savings/}{CourtListener}; \href{https://www.law.cornell.edu/supct/html/98-531.ZO.html}{LII Majority}).

\begin{quote}
``Because of this lack, the provisions of the Patent Remedy Act are `so out of proportion to a supposed remedial or preventive object that [they] cannot be understood as responsive to, or designed to prevent, unconstitutional behavior.' ''
\end{quote}
\noindent --- Chief Justice William H. Rehnquist, Opinion of the Court (quoting \textit{City of Boerne}, 521 U.S. at 532), 527 U.S. at 646 (\href{https://www.courtlistener.com/opinion/118316/florida-prepaid-postsecondary-education-expense-board-v-college-savings/}{CourtListener}; \href{https://www.law.cornell.edu/supct/html/98-531.ZO.html}{LII Majority}).

\begin{quote}
``The Court's opinion today threatens to read Congress' power to pass prophylactic legislation out of \S~5 altogether; its holding is unsupported by \textit{City of Boerne} and in fact conflicts with our reasoning in that case.''
\end{quote}
\noindent --- Justice John Paul Stevens, Dissenting Opinion, 527 U.S. at 660 (\href{https://www.courtlistener.com/opinion/118316/florida-prepaid-postsecondary-education-expense-board-v-college-savings/}{CourtListener}; \href{https://www.law.cornell.edu/supct/html/98-531.ZD.html}{LII Dissent}).


\end{document}
